\begin{solapartat}
A partir de $\lambda = 400\ \mathrm{nm}$ obtenim la freqüència dels fotons incidents
\[ \lambda = \frac{c}{f} \ \text{\punts{0,1}} \quad\to\quad f = \frac{c}{\lambda} = \frac{3\cdot 10^{8}}{4\cdot 10^{-7}} = 7,50\cdot 10^{14}\ \mathrm{Hz} \ \text{\punts{0,2}} \]
i la seva energia
\[ h\,f \ \text{\punts{0,1}} = (6,63\cdot 10^{-34})(7,5\cdot 10^{14}) \equiv 4,97\cdot 10^{-19}\ \mathrm{J} \ \text{\punts{0,2}}. \]
El treball d'extracció del Potassi és
\[ 2,29\ \mathrm{eV}\ \frac{1,60\cdot 10^{-19}\ \mathrm{J}}{1\ \mathrm{eV}} = 3,66\cdot 10^{-19}\ \mathrm{J}\ . \]
i la energia cinètica amb la que surten els electrons arrancats de l'elèctrode A és per tant
\[ E_c^A = h\,f - W_0 = 4,97\cdot 10^{-19} - 3,66\cdot 10^{-19} = 1,31\cdot 10^{-19}\ \mathrm{J} \ \text{\punts{0,2}}. \]
Finalment, com que $E_c = mv^2/2$, obtenim
\[ v = \sqrt{\frac{2\,E_c}{m}} = \sqrt{\frac{2\,(1,31\cdot 10^{-19})}{9,11\cdot 10^{-31}}} = 5,36\cdot 10^{5}\ \mathrm{m/s} \ \text{\punts{0,2}}. \]
\end{solapartat}

\begin{solapartat}
El potencial de frenada és el valor mínim de tensió que fa que els electrons que surten d'un dels elèctrodes no arribin a l'altre. Per tal d'aconseguir això, l'elèctrode B ha d'estar a un potencial menor que l'elèctrode A. Així doncs, quan incideixen fotons de freqüència $f$ sobre A, s'arranquen electrons amb energia cinètica $E_c^A$ d'acord a l'expressió
\[ h\,f = W_0 + E_c^A\ , \ \text{\punts{0,2}} \]
mentre que l'equació del balanç d'energia dels electrons que surten de A i van a B ens diu que
\[ E_c^A + E_p^A = E_c^B + E_p^B \ \text{\punts{0,2}} \]
sent $E_c$ i $E_p$ les energies cinètica i potencial elèctrica, respectivament. En el nostre cas haurà de ser $E_c^B = 0$ i per tant
\[ E_c^A = E_p^B - E_p^A = q(V_B - V_A) \equiv -q\,V_f \ \text{\punts{0,1}} \]
on, al ser $q < 0$, veiem que $V_A > V_B$ com era d'esperar. A l'anterior expressió, $V_f$ és el valor del potencial de frenada. Juntant les expressions anteriors trobem
\[ h\,f - W_0 = -q\,V_f \ \text{\punts{0,1}} \]
i per tant
\[ W_0 = h\,f + q\,V_f\ . \]
A partir dels resultat d'abans per a la llum de longitud d'ona $\lambda = 400\ \mathrm{nm}$ obtinguts abans
\[ h\,f = 4,97\cdot 10^{-19}\mathrm{J}\ \frac{1\ \mathrm{eV}}{1,60\cdot 10^{-19}\ \mathrm{J}} = 3,11\ \mathrm{eV}\ . \]
obtenim
\[ W_0 = h\,f + q\,V_f = 3,11 - 0,17 = 2,94\ \mathrm{eV} \ \text{\punts{0,2}} \]
de forma que el material desconegut és Liti. \punts{0,2}
\end{solapartat}
